\documentclass[11pt]{article}
\usepackage{ap-psych}
\title{Unit 4: Social Psychology and Personality}
\author{Paul Wang}
\date{November 2026}
\docpart{Unit 4 \textperiodcentered\ Study Guide \textperiodcentered\ 15--25\% of MCQ}

\begin{document}
\maketitle

\begin{abstract}
Unit 4 explains why people think and act as they do in relation to others, what stable traits they carry from situation to situation, and what drives and colors their behavior. It covers attribution and person perception (4.1), attitudes and persuasion (4.2), conformity, obedience, and group behavior (4.3), psychodynamic and humanistic personality theories (4.4), social-cognitive and trait theories (4.5), motivation (4.6), and emotion (4.7). Questions here are almost always scenario-based and the distractors are near-synonyms, so the tables pair the terms most often tested against each other. Each section ends with check questions; the key is on the last page.
\end{abstract}

\scoringline{Check questions: \blank\ / \qtotal \qquad Target: 20+ \qquad Bias pairs to drill: \blank[10em]}

\section{Attribution and person perception (4.1)}

An \textbf{attribution} explains behavior by the person (\textbf{dispositional}: traits, effort) or the situation (\textbf{situational}: luck, pressure, circumstance).
\begin{terms}
Fundamental attribution error & Overestimating dispositional and underestimating situational causes of \emph{others'} behavior (``He's late because he's lazy''). Stronger in individualist cultures. \\
Actor--observer bias & Explaining your own behavior situationally but others' dispositionally. \\
Self-serving bias & Taking credit for successes (dispositional) and blaming failures on the situation. \\
Explanatory style & \textbf{Optimistic}: bad events seen as external, temporary, specific. \textbf{Pessimistic}: internal, stable, global. \\
Locus of control & \textbf{Internal}: I control my outcomes. \textbf{External}: outside forces do. \\
Mere-exposure effect & Repeated exposure increases liking. \\
Self-fulfilling prophecy & An expectation leads to behavior that makes it come true. \\
Social comparison & Judging yourself against others. \textbf{Upward} (to those better off) can motivate or discourage; \textbf{downward} can boost self-esteem. \textbf{Relative deprivation}: feeling worse off than those you compare yourself to. \\
\end{terms}

\begin{q}{2}{4.1 Attribution}{Blaming a stranger's traits and ignoring the situation = fundamental attribution error.}
A driver cuts Jen off and she thinks, ``What a reckless jerk,'' without considering he might be rushing to a hospital. Jen is committing
\ch{the self-serving bias}{the fundamental attribution error}{the just-world phenomenon}{relative deprivation}
\end{q}

\begin{q}{4}{4.1 Self-serving bias}{Credit for success to self, blame for failure to situation.}
Marco says his A in chemistry came from his hard work, but his D in history was because the teacher was unfair. This is
\ch{the fundamental attribution error}{the mere-exposure effect}{an external locus of control}{the self-serving bias}
\end{q}

\begin{q}{1}{4.1 Self-fulfilling prophecy}{Teacher expectations change teacher behavior, producing the expected gains.}
Teachers told that certain (randomly chosen) students were ``late bloomers'' gave them more attention, and those students' scores rose. This illustrates
\chh{a self-fulfilling prophecy}{social loafing}{the halo effect}{downward comparison}
\end{q}

\section{Attitudes and persuasion (4.2)}

\begin{terms}
Stereotype & An overgeneralized \emph{belief} about a group. \\
Prejudice & An unjustified (usually negative) \emph{attitude} toward a group. \\
Discrimination & Unjustified negative \emph{behavior} toward group members. \\
Implicit attitudes & Automatic associations people may not know they hold (measured by the Implicit Association Test). \\
In-group bias & Favoring one's own group. \textbf{Ethnocentrism}: judging other cultures by one's own standards. \\
Out-group homogeneity & Seeing members of other groups as all alike (``they're all the same''). \\
Just-world phenomenon & Believing people get what they deserve, which can lead to blaming victims. \\
Belief perseverance & Clinging to beliefs even after the evidence for them is discredited. \\
Cognitive dissonance & Discomfort when actions and attitudes conflict; people often change the \emph{attitude} to match the action (Festinger). \\
Halo effect & One positive trait colors overall impressions. \\
\end{terms}

\textbf{Elaboration likelihood model}: the \textbf{central route} persuades with strong arguments and evidence when people are motivated and able to think (lasting change); the \textbf{peripheral route} uses cues such as attractiveness, celebrity, or repetition (quick, shallow change). \textbf{Foot-in-the-door}: a small request first increases compliance with a larger one. \textbf{Door-in-the-face}: an outrageous request first makes a smaller one seem reasonable.

\begin{q}{3}{4.2 Dissonance}{Low pay for a dull task \(\rightarrow\) attitude shifts to justify the lie (Festinger).}
Participants paid only \$1 to tell others a boring task was fun later rated the task as more enjoyable than those paid \$20. This is best explained by
\chh{the halo effect}{belief perseverance}{cognitive dissonance}{the peripheral route}
\end{q}

\begin{q}{2}{4.2 Persuasion}{Large request first, then the real (smaller) request = door-in-the-face.}
A charity first asks you to volunteer every weekend for a year. When you refuse, it asks you to help for one afternoon, and you agree. This is the
\ch{foot-in-the-door technique}{door-in-the-face technique}{central route to persuasion}{mere-exposure effect}
\end{q}

\begin{q}{4}{4.2 ELM}{Celebrity endorsement without arguments = peripheral route.}
An ad features a famous athlete smiling with a sports drink but gives no information about the product. It relies on
\ch{the central route to persuasion}{cognitive dissonance}{belief perseverance}{the peripheral route to persuasion}
\end{q}

\begin{q}{1}{4.2 Just-world}{Assuming victims deserve their fate = just-world phenomenon.}
After a woman is robbed, a neighbor says, ``She must have done something to invite it.'' This reflects
\chh{the just-world phenomenon}{out-group homogeneity}{the halo effect}{ethnocentrism}
\end{q}

\section{Social situations and groups (4.3)}

\subsection{Conformity and obedience}
\textbf{Asch}: about one-third of responses conformed to an obviously wrong group answer about line lengths. Conformity rises with a unanimous, larger (up to about 3--4), high-status group, public responses, and collectivist cultures. \textbf{Normative social influence}: conforming to be accepted; \textbf{informational social influence}: conforming because others seem to know better.

\textbf{Milgram}: 65\% of participants delivered what they believed was the maximum 450-volt shock when ordered by an authority. Obedience was highest when the authority was close and legitimate and the victim was distant; it fell when others defied the order. The study is also a classic case of ethical controversy (stress, deception).

\subsection{Group effects}
\begin{terms}
Social facilitation & Better performance on simple or well-learned tasks in front of others; worse on hard ones. \\
Social loafing & Less effort in a group when individual work is not tracked. \\
Deindividuation & Loss of self-awareness and restraint in anonymous, arousing groups. \\
Group polarization & Group discussion strengthens members' prevailing views. \\
Groupthink & Desire for harmony overrides realistic evaluation of options. \\
Diffusion of responsibility & Each person feels less responsible when others are present. \\
Bystander effect & Any one person is less likely to help when others are present. \\
False consensus effect & Overestimating how much others share your beliefs. \\
Social traps & Pursuing self-interest leads to collective harm (overfishing). \\
Superordinate goals & Shared goals that require cooperation; reduce intergroup conflict (Sherif's Robbers Cave). \\
\end{terms}
\textbf{Individualist} cultures prioritize personal goals and identity; \textbf{collectivist} cultures prioritize group goals and identity. \textbf{Prosocial behavior} is guided by the \textbf{reciprocity norm} (help those who help you) and the \textbf{social responsibility norm} (help those who need it). \textbf{Altruism} is unselfish concern for others. \textbf{Industrial-organizational psychology} applies these ideas to the workplace.

\begin{q}{3}{4.3 Bystander}{More witnesses \(\rightarrow\) each feels less responsible: diffusion of responsibility.}
A student collapses in a crowded hallway and dozens walk past assuming someone else will help. This is best explained by
\chh{groupthink}{social facilitation}{diffusion of responsibility}{the reciprocity norm}
\end{q}

\begin{q}{1}{4.3 Group polarization}{Like-minded discussion intensifies the shared view.}
Students who mildly support a later school start become much more strongly in favor after discussing it only with each other. This is
\chh{group polarization}{groupthink}{deindividuation}{social loafing}
\end{q}

\begin{q}{4}{4.3 Social facilitation}{Others' presence improves well-learned performance and worsens new, hard tasks.}
An expert pianist plays better in front of an audience, while a beginner plays worse. This pattern is
\chh{social loafing}{conformity}{deindividuation}{social facilitation}
\end{q}

\begin{q}{2}{4.3 Conformity}{Going along with others because they seem to know = informational influence.}
In a foreign city, Tom follows the crowd off the train because he assumes they know it is the right stop. This is
\ch{normative social influence}{informational social influence}{obedience to authority}{the false consensus effect}
\end{q}

\section{Psychodynamic and humanistic theories (4.4)}

\subsection{Psychodynamic}
Behavior is shaped by unconscious motives and early childhood experience. Freud's structure: the \textbf{id} (pleasure principle, immediate gratification), \textbf{superego} (internalized morals; guilt), and \textbf{ego} (reality principle; mediates). Anxiety arises from their conflict; the ego defends with \textbf{defense mechanisms}:
\begin{terms}
Denial & Refusing to accept reality (``I'm not addicted''). \\
Displacement & Shifting impulses to a safer target (yelling at your sibling after your boss yells at you). \\
Projection & Attributing your own unacceptable feelings to others. \\
Rationalization & Offering a self-justifying but false explanation. \\
Reaction formation & Expressing the opposite of a true feeling. \\
Regression & Retreating to an earlier developmental stage's behavior. \\
Repression & Pushing threatening thoughts and memories out of awareness. \\
Sublimation & Channeling unacceptable impulses into socially valued activities (aggression into boxing). \\
\end{terms}
\textbf{Projective tests} present ambiguous stimuli to reveal unconscious content: \textbf{Rorschach inkblots}, \textbf{Thematic Apperception Test} (TAT). Critics note low reliability and validity.

\subsection{Humanistic}
Emphasizes free will and growth. \textbf{Maslow}: \textbf{self-actualization}, fulfilling one's potential, at the top of a hierarchy of needs. \textbf{Rogers}: growth requires \textbf{unconditional positive regard} (acceptance without conditions), genuineness, and empathy. Healthy personality depends on \textbf{congruence} between one's \textbf{self-concept} and experience.

\begin{q}{3}{4.4 Defense mechanisms}{Redirecting anger to a safer target = displacement.}
After being scolded by his coach, Dev goes home and kicks his dog. Freud would call this
\chh{projection}{sublimation}{displacement}{reaction formation}
\end{q}

\begin{q}{1}{4.4 Defense mechanisms}{Acting the opposite of a forbidden feeling = reaction formation.}
A man who is attracted to his best friend's partner is excessively rude to her. This reflects
\chh{reaction formation}{regression}{rationalization}{denial}
\end{q}

\begin{q}{2}{4.4 Humanistic}{Acceptance without conditions = Rogers's unconditional positive regard.}
A parent tells a child, ``I love you no matter what grades you get.'' Rogers would say the parent is offering
\ch{conditional positive regard}{unconditional positive regard}{a projective test}{a superego demand}
\end{q}

\section{Social-cognitive and trait theories (4.5)}

\textbf{Social-cognitive} (Bandura): personality comes from interaction of traits, thoughts, and environment. \textbf{Reciprocal determinism}: behavior, personal factors, and environment influence one another. \textbf{Self-efficacy}: belief in one's ability to succeed at a specific task. \textbf{Self-concept}, \textbf{self-esteem} (overall sense of self-worth).

\textbf{Trait theories} describe stable, measurable dispositions, identified through \textbf{factor analysis}. The \textbf{Big Five}: \textbf{O}penness, \textbf{C}onscientiousness, \textbf{E}xtraversion, \textbf{A}greeableness, \textbf{N}euroticism (emotional stability). Traits are fairly stable in adulthood and partly heritable. \textbf{Personality inventories} (long self-report questionnaires) measure traits; they are more reliable and valid than projective tests.

\begin{q}{4}{4.5 Self-efficacy}{Confidence in one's ability to succeed at a specific task = self-efficacy.}
Despite failing her first try, Ana is confident she can pass her driving test because she has mastered each skill. This is high
\chh{self-esteem}{neuroticism}{external locus of control}{self-efficacy}
\end{q}

\begin{q}{3}{4.5 Big Five}{Organized, dependable, disciplined = conscientiousness.}
A student who is organized, plans ahead, and always meets deadlines would score high in
\chh{openness}{agreeableness}{conscientiousness}{extraversion}
\end{q}

\begin{q}{1}{4.5 Reciprocal determinism}{Traits, behavior, and environment influence each other: reciprocal determinism.}
An outgoing student joins drama club, the club's friendships make her even more outgoing, and she seeks more performances. This cycle best illustrates
\ch{reciprocal determinism}{the id--superego conflict}{unconditional positive regard}{a projective test}
\end{q}

\section{Motivation (4.6)}

\begin{compare}{Theory}{Core idea}{Example}
Instinct & Unlearned, species-wide behaviors drive action. & A spider weaves a web. \\
Drive-reduction & Needs create aversive drives that push us to restore \textbf{homeostasis}. & Thirst drives drinking. \\
Arousal & We seek an optimal level of arousal. \textbf{Yerkes--Dodson law}: performance peaks at moderate arousal (lower for hard tasks, higher for easy ones). & Mild nerves help on a test; panic hurts. \\
Incentive & External rewards pull behavior. & Working for a bonus. \\
Self-determination & Motivation thrives when needs for \textbf{autonomy}, \textbf{competence}, and \textbf{relatedness} are met. & Choosing your own project. \\
\end{compare}
\textbf{Intrinsic motivation}: doing something for its own sake. \textbf{Extrinsic}: for external rewards. The \textbf{overjustification effect}: rewarding an already-enjoyed activity can reduce intrinsic motivation. \textbf{Sensation seeking}: the drive for novel, intense experiences. \textbf{Maslow's hierarchy}: physiological \(\rightarrow\) safety \(\rightarrow\) belonging \(\rightarrow\) esteem \(\rightarrow\) self-actualization.

\textbf{Conflicts}: \textbf{approach--approach} (two good options), \textbf{avoidance--avoidance} (two bad options), \textbf{approach--avoidance} (one option with good and bad parts), \textbf{multiple approach--avoidance} (several options, each mixed).

\textbf{Hunger}: the \textbf{hypothalamus} regulates; \textbf{ghrelin} signals hunger, \textbf{leptin} signals fullness; set point and culture also shape eating.

\begin{q}{2}{4.6 Yerkes-Dodson}{Performance on hard tasks peaks at lower arousal.}
According to the Yerkes--Dodson law, a student taking a very difficult exam will perform best with
\chh{very high arousal}{relatively low-to-moderate arousal}{no arousal at all}{maximum anxiety}
\end{q}

\begin{q}{4}{4.6 Overjustification}{Paying for an enjoyed activity can undermine intrinsic motivation.}
A child who loves drawing is paid for each picture; when the payments stop, she draws less than before. This is
\chh{drive reduction}{sensation seeking}{an approach--avoidance conflict}{the overjustification effect}
\end{q}

\begin{q}{1}{4.6 Conflicts}{One option with both attractive and unattractive features = approach-avoidance.}
Lia wants the high-paying job but dreads its two-hour commute. She faces
\ch{an approach--avoidance conflict}{an approach--approach conflict}{an avoidance--avoidance conflict}{homeostasis}
\end{q}

\section{Emotion (4.7)}

\begin{compare}{Theory}{Sequence}{Key claim}
James--Lange & Stimulus \(\rightarrow\) physiological arousal \(\rightarrow\) emotion & ``We feel afraid because we tremble.'' \\
Cannon--Bard & Stimulus \(\rightarrow\) arousal and emotion \emph{simultaneously} & The thalamus sends signals to body and cortex at once. \\
Schachter--Singer two-factor & Arousal + cognitive label \(\rightarrow\) emotion & The same arousal can be felt as fear or excitement depending on how it is interpreted. \\
\end{compare}
\textbf{Facial-feedback hypothesis}: facial expressions intensify the emotions they show (holding a pen in your teeth to force a smile). Basic emotional expressions (happiness, sadness, fear, anger, surprise, disgust) are recognized \textbf{universally} (Ekman), but \textbf{display rules}, cultural norms for when to show emotion, vary. \textbf{Broaden-and-build theory} (Fredrickson): positive emotions widen thinking and build lasting resources.

\begin{q}{3}{4.7 Two-factor}{Arousal labeled by context (the excited crowd) = Schachter-Singer two-factor.}
Fans in a stadium feel their hearts pounding and, seeing everyone cheering, interpret the feeling as excitement. This fits
\ch{the James--Lange theory}{the Cannon--Bard theory}{the Schachter--Singer two-factor theory}{the facial-feedback hypothesis}
\end{q}

\begin{q}{2}{4.7 James-Lange}{Body response first, then the emotion = James-Lange.}
``I see a snake, my heart races, and I recognize from my racing heart that I am afraid.'' This sequence matches
\ch{the Cannon--Bard theory}{the James--Lange theory}{the two-factor theory}{broaden-and-build theory}
\end{q}

\begin{q}{4}{4.7 Display rules}{Cultural norms for showing emotion = display rules.}
Japanese participants smiled to mask disgust when an authority figure was present but showed disgust when alone. This demonstrates
\chh{universal expressions}{the facial-feedback hypothesis}{Cannon--Bard theory}{display rules}
\end{q}

\makekey
\end{document}
